\documentclass[oneside,a4paper,14pt]{extarticle} %размер шрифта 14
\usepackage[T1,T2A]{fontenc}
\usepackage[
a4paper,
letterpaper,
top=2cm,
bottom=2cm,
left=2.5cm,
right=1.5cm
]{geometry}
\usepackage[utf8]{inputenc} %кодировка текста
\usepackage[russian]{babel} %поддержка русского языка
\usepackage{textcomp} %текстовые символы
\usepackage{indentfirst} %корректировка отступов
\usepackage{listings} %для кода !!!
\usepackage{listingsutf8}
\usepackage{xcolor}
\usepackage{float}
\usepackage{enumitem} %для нумерации буквами
\usepackage{graphicx} %работа с изображениям
\usepackage{mwe} % for blindtext and example-image-a in example
\usepackage{wrapfig}
\usepackage{caption}
\usepackage{amsmath}  % для формул и символов
\usepackage{amsfonts}
\usepackage{amsthm}
\usepackage[breaklinks]{hyperref}
\usepackage{xurl}
% размер кегля у заголоков разделов
\usepackage{titlesec}
\titleformat{\section}
{\normalfont\large\bfseries}{\thesection}{1em}{}
\titleformat{\subsection}
{\normalfont\normalsize\bfseries}{\thesubsection}{1em}{}
\titleformat{\subsubsection}
{\normalfont\normalsize\bfseries}{\thesubsection}{1em}{}
\renewcommand\baselinestretch{1.33}
\normalsize % межстрочный интервал
\setlength{\parindent}{1.25cm}
% настройки листинга
\lstset{
language={},
basicstyle=\ttfamily\small,
keywordstyle=\color{blue},
commentstyle=\color{green!60!black},
stringstyle=\color{red},
numbers=left,
numberstyle=\tiny\color{gray},
frame=single,
breaklines=true,
showstringspaces=false,
tabsize=2,
inputencoding=cp1251,
extendedchars=true,
}
\begin{document}
\newpage
\thispagestyle{empty}
\begin{center}
МИНИСТЕРСТВО НАУКИ И ВЫСШЕГО ОБРАЗОВАНИЯ\\
РОССИЙСКОЙ ФЕДЕРАЦИИ
ФЕДЕРАЛЬНОЕ ГОСУДАРСТВЕННОЕ БЮДЖЕТНОЕ\\
ОБРАЗОВАТЕЛЬНОЕ
УЧРЕЖДЕНИЕ ВЫСШЕГО ОБРАЗОВАНИЯ\\
«ВЯТСКИЙ ГОСУДАРСТВЕННЫЙ УНИВЕРСИТЕТ»\\
Институт математики и информационных систем\\
Факультет автоматики и вычислительной техники\\
Кафедра электронных вычислительных машин
\end{center}
\vspace{20mm}
\begin{center}
Отчёт по лабораторной работе №8\\
по дисциплине\\
<<Промпт-инжиниринг>>\\
<<Создание синтетических наборов данных для машинного обучения>>\\
<<Вариант 4>>\\
\end{center}
\vspace{48mm}
\noindent
Выполнил студент гр. ИВТб-1303-06-00 \hfill \rule[-0,5mm]{30mm}{0.15mm}\,/Гортоломей И.К./
\newline
Проверил доцент кафедры ЭВМ \hfill \rule[-0,5mm]{30mm}{0.15mm}\,/Колупаев А.В./
\vfill
\begin{center}
Киров\\
2026
\end{center}
\newpage
\thispagestyle{empty}
\section*{Введение}
\noindent
\textbf{Цель работы:} Формирование практических навыков проектирования, генерации и валидации синтетических наборов данных с использованием больших языковых моделей. Освоение полного цикла подготовки данных: от формализации схемы и генерации табличных структур (CSV) до создания семантически размеченных корпусов для NLP-задач и автоматизации процессов контроля качества (Data Quality Assurance) без написания программного кода.
\noindent
\textbf{Используемые модели и инструменты:}
\begin{itemize}
\item \textbf{DeepSeek} --- основная языковая модель для всех трёх заданий.
\item \textbf{GigaChat} --- вторая модель для сравнительного анализа.
\item \textbf{Google Sheets} --- табличный процессор для валидации CSV-файлов.
\end{itemize}
\noindent
\textbf{Предметная область варианта 4:}
\begin{itemize}
\item \textbf{Задание 1 и 3:} Логистика --- отслеживание почтовых отправлений (система трекинга посылок).
\item \textbf{Задание 2:} Intent Recognition --- голосовой помощник банка (определение намерений пользователя).
\end{itemize}
\newpage
\section{Задание №1: Генерация структурированного набора данных}
\subsection{Постановка задачи варианта}
\noindent
\textbf{Предметная область:} Логистика --- отслеживание почтовых отправлений.
\noindent
\textbf{Атрибуты:}
\begin{itemize}
\item Трек-номер (строка, формат TRK-XXXXXX)
\item Вес (кг) (число с плавающей точкой, 0.1--30.0)
\item Город отправления (строка)
\item Город назначения (строка)
\item Дата отправки (ГГГГ-ММ-ДД)
\item Текущий статус (В пути/Доставлено/Таможня)
\item Флаг хрупкого груза (Да/Нет)
\end{itemize}
\noindent
Требуется сгенерировать CSV-файл с 20 строками, соблюдая типы данных и логические ограничения (вес положительный, дата в прошлом или настоящем, статус из заданного списка).
\subsection{Гипотеза и проектирование промпта}
\noindent
Первоначальный промпт (нулевая версия) содержал лишь описание задачи без примеров формата:
\lstinputlisting[language=]{1.nonutf.txt}
\newpage
\noindent
Результат содержал пояснения, нарушение формата весов (смесь кг и г), разные форматы дат и пропуски столбцов. Для исправления применены:
\begin{itemize}
\item \textbf{Ролевая установка} --- модель как генератор тестовых данных логистической системы.
\item \textbf{Инструкции по форматированию} --- выводить только CSV, без лишнего текста.
\item \textbf{Few-Shot} --- заголовок и 2 примера строк.
\item \textbf{Негативные ограничения} --- запрет пояснений, нумерации, лишних кавычек.
\item \textbf{Явные правила проверки} --- вес > 0, статус из списка, флаг Да/Нет.
\end{itemize}
\subsection{Журнал итеративной отладки}
\noindent
\textbf{Итерация 1} \\
\noindent
\textbf{Вход:} добавлена ролевая установка и требование выводить только CSV. \\
\noindent
\textbf{Выход (DeepSeek):}
\lstinputlisting[language=]{2.nonutf.txt}
\noindent
\textbf{Анализ дефектов:} В столбце «Вес» встречались отрицательные значения и некорректные форматы. \\
\noindent
\textbf{Корректирующее действие:} Добавлено правило: «Вес --- положительное число от 0.1 до 30.0 кг».
\newpage
\noindent
\textbf{Итерация 2} \\
\noindent
\textbf{Вход:} добавлено правило корректного веса и проверка статусов. \\
\noindent
\textbf{Выход (GigaChat):}
\lstinputlisting[language=]{3.nonutf.txt}
\noindent
\textbf{Анализ дефектов:} Наличие пояснительного текста «Вот ваш CSV:» мешает автоматическому парсингу. \\
\noindent
\textbf{Корректирующее действие:} Усилено негативное ограничение: «Не выводи никаких пояснений, вступлений или заключений. Начинай сразу с заголовка». \\
\newline
\noindent
\textbf{Итерация 3 (финальная, сравнение моделей)} \\
\noindent
\textbf{Вход:} добавлен Few-Shot пример с двумя строками, демонстрирующими корректный формат. \\
\noindent
\textbf{Выход (DeepSeek):}
\lstinputlisting[language=]{4.nonutf.txt}
\noindent
\textbf{Выход (GigaChat) на том же финальном промпте:}
\lstinputlisting[language=]{5.nonutf.txt}
\noindent
\textbf{Анализ дефектов:} Обе модели справились хорошо. GigaChat иногда добавлял лишние пробелы в названиях городов. \\
\noindent
\textbf{Корректирующее действие:} Промпт остаётся без изменений.
\newpage
\subsection{Демонстрация валидации результата}
\noindent
\begin{figure}[H]
\centering
\includegraphics[width=17cm]{1.png}
\caption{Демонстрация валидации CSV-файла в табличном редакторе}
\end{figure}
\subsection{Сравнительный анализ моделей (Benchmarking)}
\noindent
\begin{table}[H]
\centering
\begin{tabular}{|p{4cm}|p{6cm}|p{6cm}|}
\hline
\textbf{Критерий} & \textbf{DeepSeek} & \textbf{GigaChat} \\ \hline
Синтаксическая корректность CSV & Идеально: ни одного лишнего слова, ровно 20 строк & После первой попытки добавил пояснение, после уточнения --- нормально \\ \hline
Следование схеме данных & Все 7 столбцов, правильные типы, веса в диапазоне 0.1-30.0 & Все столбцы, но в некоторых строках лишние пробелы в городах \\ \hline
Точность следования ограничениям & 98\% & 92\% \\ \hline
Итоговая оценка пригодности & Готов к использованию без правок & Требуется незначительная проверка форматов \\ \hline
\end{tabular}
\end{table}
\noindent
\textbf{Обоснование:} DeepSeek строже соблюдает форматы и меньше «фантазирует» с пояснениями. GigaChat на финальном промпте работал хорошо, но требовал более тщательной настройки негативных ограничений.
\newpage
\subsection{Финальный промпт (задание 1)}
\lstinputlisting[language=]{6.nonutf.txt}
\noindent
(В работе использовано \texttt{КОЛИЧЕСТВО\_СТРОК} = 20)
\newpage
\section{Задание №2: Генерация размеченного набора данных}
\subsection{Постановка задачи варианта}
\noindent
\textbf{Предметная область:} Intent Recognition --- Голосовой помощник банка. \\
\noindent
\textbf{Классы разметки (намерения):}
\begin{enumerate}
\item Узнать баланс
\item Заблокировать карту
\item Перевод средств
\item Оформить кредит
\item Связаться с оператором
\end{enumerate}
\noindent
Требуется сгенерировать 30 фраз (по 6 на класс) с использованием техники персон для лексического разнообразия.
\subsection{Гипотеза и проектирование промпта}
\noindent
Начальная версия промпта задавала только классы и просила сгенерировать по 5 примеров на класс. Результаты были однообразны. Применены:
\begin{itemize}
\item \textbf{Техника персон} --- 6 профилей авторов.
\item \textbf{Few-Shot} --- показаны примеры формата CSV с двумя колонками.
\item \textbf{Явное требование разнообразия формулировок}.
\item \textbf{Обязательные кавычки вокруг текста} для корректного парсинга CSV.
\end{itemize}
\newpage
\subsection{Журнал итеративной отладки}
\noindent
\textbf{Итерация 1} \\
\noindent
\textbf{Вход:} простой запрос: «Сгенерируй 5 фраз для каждого намерения пользователя в банке» \\
\noindent
\textbf{Выход (DeepSeek):} 
\lstinputlisting[language=]{7.nonutf.txt}
\noindent
\textbf{Анализ дефектов:} отсутствие имитации реальных пользователей, лексическая бедность, повторы конструкций. \\
\noindent
\textbf{Корректирующее действие:} добавлены профили персон и требование использовать разные эмоциональные окраски и стили речи. \\
\newline
\noindent
\textbf{Итерация 2} \\
\noindent
\textbf{Вход:} добавлены 6 профилей, но без привязки к конкретным классам \\
\noindent
\textbf{Выход (GigaChat):}
\lstinputlisting[language=]{8.nonutf.txt}
\noindent
\textbf{Анализ дефектов:} для некоторых классов не было примеров от всех типов пользователей; стили не всегда соответствовали персоне. \\
\noindent
\textbf{Корректирующее действие:} явно указано: «Для каждого сочетания (класс + профиль) сгенерируй по одной фразе»; добавлено требование строгого соответствия стиля персоне.
\newpage
\noindent
\textbf{Итерация 3 (финальная, сравнение моделей)} \\
\noindent
\textbf{Вход:} финальный промпт с перечислением 5 классов и 6 профилей, требование выводить CSV с кавычками \\
\noindent
\textbf{Выход (DeepSeek):} 30 строк без ошибок (фрагмент):
\lstinputlisting[language=]{9.nonutf.txt}
\noindent
\textbf{Выход (GigaChat) на том же финальном промпте:}
\lstinputlisting[language=]{10.nonutf.txt}
\noindent
\textbf{Анализ дефектов:} GigaChat допустил незначительные отклонения в формулировках меток. \\
\noindent
\textbf{Корректирующее действие:} Промпт остаётся без изменений (рекомендуется использовать DeepSeek для высокого качества).
\newpage
\subsection{Демонстрация валидации результата}
\noindent
Ниже представлен фрагмент (первые 10 строк) итогового CSV-файла, сгенерированного DeepSeek:
\lstinputlisting[language=]{11.nonutf.txt}
\subsection{Сравнительный анализ моделей (Benchmarking)}
\noindent
\begin{table}[H]
\centering
\begin{tabular}{|p{6cm}|p{5cm}|p{5cm}|}
\hline
\textbf{Критерий} & \textbf{DeepSeek} & \textbf{GigaChat} \\ \hline
Семантическая согласованность & 100\% & 93\% (2 ошибки в названиях меток) \\ \hline
Соблюдение ролевой модели & Ярко выражены все 6 профилей & Рассерженный выражен слабо \\ \hline
Формат CSV & Безупречно & Иногда добавлял пояснения \\ \hline
Итоговая оценка пригодности & Готов к использованию & Требуется ручная проверка \\ \hline
\end{tabular}
\end{table}
\noindent
\textbf{Обоснование:} DeepSeek лучше удерживает требование разнообразия и точнее привязывает стиль к персоне. GigaChat генерирует более «сглаженные» фразы, но для стресс-тестирования классификатора это недостаточно.
\subsection{Финальный промпт (задание 2)}
\lstinputlisting[language=]{12.nonutf.txt}
\newpage
\section{Задание №3: Контроль качества и очистка данных}
\subsection{Постановка задачи варианта}
\noindent
Используя предметную область из Задания 1 (логистика --- отслеживание посылок), требуется:
\begin{enumerate}
\item Сгенерировать «грязный» набор данных с внедрением аномалий:
\begin{itemize}
\item нарушение формата трек-номера (TRK-1, TRK001, без префикса)
\item нечисловые значения в весе («легкая», «N/A», отрицательные)
\item опечатки в названиях городов («Мосвка», «Санкт-Петербруг»)
\item некорректные статусы («впути», «Доставлен», «Таможня clearance»)
\item пропуски полей (пустые ячейки)
\item логические противоречия (флаг хрупкого груза для веса 30 кг)
\end{itemize}
\item Разработать промпт-очиститель, который исправляет все ошибки в соответствии с правилами, не изменяя корректные данные.
\end{enumerate}
\subsection{Гипотеза и проектирование промпта}
\noindent
Для генерации «грязных» данных использована техника \textbf{Constraint Injection} --- в промпт включены явные инструкции по внедрению ошибок с указанием частоты (20--30\% строк). Промпт-очиститель построен как набор формальных правил с требованием \textbf{Chain-of-Thought} и строгим запретом на галлюцинации.
\newpage
\subsection{Журнал итеративной отладки}
\noindent
\textbf{Итерация 1 (генерация грязных данных)} \\
\noindent
\textbf{Вход:} промпт с перечислением 5 типов ошибок и требованием внедрить их в 20--30\% строк. \\
\noindent
\textbf{Выход (DeepSeek):} CSV с 20 строками, 6 из которых содержали аномалии (пример):
\lstinputlisting[language=]{13.nonutf.txt}
\noindent
\textbf{Анализ дефектов:} ошибки разнообразны, пропуски оформлены корректно (две запятые подряд). \\
\noindent
\textbf{Корректирующее действие:} промпт оставлен без изменений, набор принят.
\newpage
\noindent
\textbf{Итерация 2 (очистка) --- неудачная} \\
\noindent
\textbf{Вход:} первая версия очистителя с общими правилами («исправь ошибки в трек-номерах, весах, городах»). \\
\noindent
\textbf{Выход (GigaChat):} модель исправила большинство ошибок, но изменила корректные данные:
\lstinputlisting[language=]{14.nonutf.txt}
\noindent
\textbf{Анализ дефектов:} ложные срабатывания (изменение корректного трек-номера, удаление веса вместо исправления формата). \\
\noindent
\textbf{Корректирующее действие:} добавлено правило «Не изменяй корректные данные. Если поле соответствует стандарту --- оставь как есть». \\
\newline
\noindent
\textbf{Итерация 3 (финальная, сравнение моделей)} \\
\noindent
\textbf{Вход:} уточнены правила для каждого поля, добавлен явный пересчёт логики. \\
\noindent
\textbf{Выход (DeepSeek):} идеально очищенный CSV (фрагмент):
\lstinputlisting[language=]{15.nonutf.txt}
\newpage
\noindent
\textbf{Выход (GigaChat) на том же финальном промпте:}
\lstinputlisting[language=]{16.nonutf.txt}
\noindent
\textbf{Анализ дефектов:} GigaChat не исправил полностью все ошибки: опечатка в городе «Мосвка» осталась, отрицательный вес и «N/A» не исправлены. DeepSeek исправил всё.
\noindent
\textbf{Корректирующее действие:} Промпт остаётся без изменений, для GigaChat требуется дополнительная настройка.
\subsection{Демонстрация валидации результата}
\noindent
Ниже приведены «грязные» данные (первые 5 строк из 20) и результат очистки (DeepSeek).
\noindent
\textbf{Грязные данные (фрагмент):}
\lstinputlisting[language=]{17.nonutf.txt}
\noindent
\textbf{Очищенные данные (DeepSeek, фрагмент):}
\lstinputlisting[language=]{18.nonutf.txt}
\subsection{Сравнительный анализ моделей (Benchmarking)}
\noindent
\begin{table}[H]
\centering
\begin{tabular}{|p{4cm}|p{6cm}|p{6cm}|}
\hline
\textbf{Критерий} & \textbf{DeepSeek} & \textbf{GigaChat} \\ \hline
Полнота очистки (Recall) & 100\% & 75\% (пропущены опечатки в городах и нечисловые веса) \\ \hline
Точность сохранения (Precision) & 100\% & 88\% (изменил некоторые корректные трек-номера) \\ \hline
Способность к логическому выводу & Высокая (правильно обработал все форматы) & Средняя (некоторые ошибки оставил) \\ \hline
Итоговая оценка пригодности & Готов к использованию без правок & Требуется доработка и ручная проверка \\ \hline
\end{tabular}
\end{table}
\noindent
\textbf{Обоснование:} DeepSeek продемонстрировал более высокую точность и полноту очистки, правильно интерпретируя правила и не внося лишних изменений. GigaChat иногда игнорировал ошибки или менял корректные значения, что снижает надёжность автоматической очистки.
\subsection{Финальный промпт (задание 3)}
\lstinputlisting[language=]{19.nonutf.txt}
\newpage
\section*{Заключение}
\noindent
\textbf{Оценка эффективности автоматизации:} Разработанные промпты полностью автоматизируют генерацию структурированных данных, размеченных NLP-корпусов и очистку «грязных» таблиц. DeepSeek продемонстрировал высокую стабильность следования форматам и логическим правилам. GigaChat требует дополнительных итераций, но также может быть использован.
\noindent
\textbf{Выявленные ограничения:}
\begin{enumerate}
\item LLM склонны к «болтливости», поэтому требуются жёсткие негативные ограничения.
\item При генерации CSV важно явно указывать разделитель и экранирование кавычек.
\item Для Intent Recognition необходимо использовать технику персон, иначе данные получаются однообразными.
\item При очистке данных критически важно запрещать модели изменять корректные поля (правило Precision).
\end{enumerate}
\noindent
\textbf{Наиболее эффективные техники промпт-инжиниринга:}
\begin{itemize}
\item \textbf{Few-Shot} --- добавление примеров значительно улучшило структуру ответов.
\item \textbf{Negative Constraints} --- запрет пояснений и лишнего текста критичен для машинной обработки.
\item \textbf{Persona Pattern} --- обеспечивает лексическое разнообразие в задании 2.
\item \textbf{Constraint Injection} --- позволяет намеренно внедрять ошибки для тестирования очистителей.
\item \textbf{Chain-of-Thought (неявный)} --- через подробные правила в промпте очистки.
\end{itemize}
\noindent
\textbf{Рекомендации по применению LLM для генерации данных:}
\begin{itemize}
\item \textbf{Для структурных данных (Задания 1 и 3):} DeepSeek показывает превосходные результаты (95-100\% точность). Требуется минимальная ручная проверка. Экономия времени: 80-90\% по сравнению с ручным созданием.
\item \textbf{Для текстовых размеченных данных (Задание 2):} DeepSeek также предпочтительнее (100\% семантическая согласованность). GigaChat требует ручной проверки 10-15\% записей. Экономия времени: 70-80\% по сравнению с ручной генерацией.
\end{itemize}
\noindent
\textbf{Вывод:} Применение больших языковых моделей для генерации синтетических данных полностью оправдано с точки зрения экономии времени и качества результата, особенно при использовании DeepSeek и правильно сконструированных промптов с применением техник Few-Shot, Persona Pattern и Negative Constraints.
\end{document}